\section{Introduction}
\label{sec:introduction}
A learned component sits within a chain of data preparation, software
construction, compilation, and execution. Explain which part of that chain
the paper studies and why the chosen boundary matters. For an AI paper,
begin with the learning problem; for an ML systems paper, begin with the
workload and deployment constraint.

\paragraph{The research gap}
Identify the input, the assumption that fails, and the resulting consequence.
A general account of learning, such as \citet{goodfellow2016}, establishes
background; a novelty claim requires comparison with specific methods.

\paragraph{The central idea}
Describe the mechanism before its implementation details. An illustrative
method might combine a task objective with a deployment constraint.
\autoref{sec:method} shows how to express that contract and its limits.

The contribution list should name concrete outputs that the paper actually
delivers:
\begin{itemize}
  \item A problem formulation with explicit assumptions and scope.
  \item A mechanism whose design choices address the stated limitation.
  \item Evidence that tests each claim and explains the remaining failures.
\end{itemize}

